% MA 214 Discussion 13 -- covers Lectures 19 and 20.
%
% Student handout. Page 1 is the concept review; the question set runs from page 2.
% Compile from inside this folder:  pdflatex Discussion13.tex

\documentclass[11pt]{article}
\usepackage[margin=1in]{geometry}
\usepackage{parskip}
\usepackage{fancyhdr}
\usepackage{array}
\usepackage{booktabs}
\usepackage{enumitem}
\usepackage{amsmath}
\usepackage{tabularx}
\usepackage{listings}

\pagestyle{fancy}
\fancyhf{}
\cfoot{\thepage}
\rhead{Discussion 13}
\lhead{MA 214: Applied Statistics}
\renewcommand{\headrulewidth}{.4mm}

\newcommand{\blankline}[1]{\underline{\hspace{#1}}}
\newcolumntype{L}{>{\raggedright\arraybackslash}X}

\lstset{
  basicstyle=\ttfamily\small,
  columns=fullflexible,
  frame=single,
  framerule=0.4pt,
  breaklines=true,
  showstringspaces=false,
  aboveskip=8pt,
  belowskip=8pt
}

% Section header: bold title over a hairline rule.
\newcommand{\parttitle}[1]{%
  \par\medskip\noindent
  {\large\bfseries #1}\par
  \vspace{-.5em}\noindent\rule{\linewidth}{0.4pt}\par\smallskip}

% Writing space.
\newcommand{\answerspace}[1]{\vspace{#1}}

% Flush-left question list: the label runs in at the margin and the body wraps
% back to the margin, so nothing is pushed far to the right.
\newlist{questions}{enumerate}{1}
\setlist[questions]{label=\textbf{Question \arabic*.}, wide=0pt, itemsep=1.4em, topsep=.8em}

\newlist{parts}{enumerate}{1}
\setlist[parts]{label=\textbf{(\Alph*)}, leftmargin=1.6em, labelsep=.45em,
  itemsep=.7em, topsep=.5em}

\begin{document}

\noindent\textbf{Name:}\blankline{2.3in}\hfill\textbf{BUID:}\blankline{1.6in}

\begin{center}
  {\Large\bfseries Discussion 13: Once You Have Samples, Everything Is Counting}
\end{center}

\noindent\textbf{Goals:} \textbf{(1)} say why a conjugate shortcut is not always available and what
grid approximation costs, \textbf{(2)} diagnose a chain from its trace plot and acceptance rate,
\textbf{(3)} use $\widehat{R}$, ESS, and MCSE for what each is for, \textbf{(4)} answer every
posterior question by counting draws, \textbf{(5)} tell a plug-in prediction from a posterior
predictive one.

\parttitle{Part 1: Concept Review}

{\small

\noindent Everything in Part 2 can be answered from this page. Keep it in front of you.

\begin{enumerate}[label=\textbf{\arabic*.}, wide=0pt, itemsep=.4em, topsep=.5em]

\item \textbf{When the shortcut runs out, and the grid.} Conjugate priors give a closed-form
posterior for only a few prior shapes; choose another, or add parameters, and
$\int f(\theta)L(\theta;x)\,d\theta$ has no formula. \textbf{Grid approximation} is the Lecture 17
hand calculation with more rows: lay a grid over the plausible range, evaluate prior $\times$
likelihood at each point, divide by the total, read summaries off the weighted grid. Two costs: a
coarse grid can get the center about right while getting a \emph{tail} probability badly wrong, and
the number of cells is (points per axis)$^{\,k}$, which collapses past two or three parameters.

\item \textbf{MCMC.} A random-walk Metropolis sampler proposes a step, computes the ratio of
posterior heights, and accepts with that probability, so it lingers where the posterior is high. The
draws are \emph{dependent} and the early ones remember the starting value, so discard a
\textbf{burn-in}.

\item \textbf{Reading a chain.} Healthy: a horizontal band of noise with no drift. \textbf{Step too
small:} nearly every proposal is accepted but the chain inches along, so the trace wanders slowly and
neighbouring draws are almost identical. \textbf{Step too large:} nearly everything is rejected, so
the trace shows long flat runs. A \emph{high acceptance rate is not a good sign}.

\item \textbf{Three numbers.} $\widehat{R}$ compares spread \emph{between} chains with spread
\emph{within} them, so near $1.00$ means the chains agree and much above about $1.01$ means they do
not. \textbf{ESS} is how many independent draws your dependent ones are worth. \textbf{MCSE}
$= \text{posterior sd}/\sqrt{\text{ESS}}$ is the simulation error in a summary, so halving it takes
four times the ESS. $\widehat{R}$ answers ``have I converged''; ESS and MCSE answer ``have I run long
enough''.

\item \textbf{With samples in hand, every question is counting.} Given draws
$\theta^{(1)},\dots,\theta^{(B)}$ from the posterior:

\begin{center}
\renewcommand{\arraystretch}{1.15}
\begin{tabular}{ll}
\toprule
\textbf{Question} & \textbf{Answer from the draws} \\
\midrule
Point estimate & their mean, or their median \\
95\% \textbf{credible interval} & their $2.5$th and $97.5$th percentiles \\
Is $\theta > c$? & the \emph{proportion} of draws above $c$ \\
Predict a new observation & draw one $y$ from the model at \emph{each} $\theta^{(i)}$ \\
\bottomrule
\end{tabular}
\end{center}

A credible interval means what students want a confidence interval to mean: given the prior and the
data, there is a $95\%$ posterior probability that $\theta$ lies in it. A Bayesian hypothesis test is
just the third row, so there is \textbf{no null distribution and no $p$-value}.

\item \textbf{Prediction: plug in, or average.} The \textbf{plug-in} prediction uses one value,
$\hat\theta$, and ignores the fact that you are unsure about it. The \textbf{posterior predictive}
draws $y$ at every posterior draw of $\theta$, so it carries both the noise in $y$ and the
uncertainty in $\theta$. It is therefore always the wider of the two.

\item \textbf{Four traps.} Treating a high acceptance rate as evidence the sampler is good. Using the
draws before discarding burn-in. Reporting the posterior sd where the MCSE was asked for. Calling a
posterior probability a $p$-value.

\end{enumerate}

}

\newpage

\parttitle{Part 2: Question Set}

\noindent\textbf{Scenario.} A whale-watching station counts sightings per night. With a
Gamma$(4,2)$ prior and $18$ sightings over $6$ nights, Lecture 18's conjugate rule gives the
posterior $\lambda \sim$ Gamma$(22, 8)$ exactly: mean $2.750$, standard deviation $0.586$, median
$2.708$, and a $95\%$ credible interval of $[\,1.723,\; 4.013\,]$, with
$P(\lambda > 3) = 0.314$. Because the exact answer is known here, it can be used to grade the
approximations.

\begin{questions}

%%%%%%%%%%%%%%%%%%%%%%%%%%%%%%%%%%%%%%
\item \textbf{Grid approximation, and what it costs.} The station's ecologist prefers a Log-Normal
prior, for which no conjugate rule exists. You approximate on a grid over $\lambda \in [0.5, 6]$.

\begin{parts}
\item State the four steps of grid approximation in order.
\answerspace{4em}

\item Applied to the \emph{known} Gamma$(22,8)$ posterior, grids of different sizes give:

\begin{center}
\renewcommand{\arraystretch}{1.2}
\begin{tabular}{lccc}
\toprule
\textbf{Grid points} & 6 & 11 & 501 \\
\midrule
Posterior mean & 2.782 & 2.750 & 2.750 \\
$P(\lambda > 3)$ & 0.155 & 0.321 & 0.313 \\
\bottomrule
\end{tabular}
\end{center}

The exact values are $2.750$ and $0.314$. Which summary does the 6-point grid get about right, and
which does it get badly wrong? Explain why a tail probability is the more fragile of the two.
\answerspace{4.5em}

\item A later model has $4$ parameters. With $100$ grid points per axis, how many cells is that?
Write the arithmetic, and say what this rules out.
\answerspace{4em}
\end{parts}

%%%%%%%%%%%%%%%%%%%%%%%%%%%%%%%%%%%%%%
\item \textbf{Tuning the sampler.} You run a random-walk Metropolis sampler on the same posterior
three times, changing only the step size. Each run is $6000$ draws with the first $1000$ discarded.

\begin{center}
\renewcommand{\arraystretch}{1.25}
\begin{tabular}{cccc}
\toprule
\textbf{Step size} & \textbf{Acceptance rate} & \textbf{ESS} & \textbf{Posterior mean from draws} \\
\midrule
0.05 & 0.987 & 14.2 & 2.678 \\
1.20 & 0.611 & 867.0 & 2.750 \\
8.00 & 0.116 & 480.5 & 2.770 \\
\bottomrule
\end{tabular}
\end{center}

\begin{parts}
\item Which run would you use, and why is it \emph{not} the one with the highest acceptance rate?
\answerspace{4em}

\item Describe what the trace plot looks like for the step of $0.05$, and separately for the step of
$8.00$. The two look nothing alike even though both are bad.
\answerspace{4.5em}

\item The step of $0.05$ accepted $98.7\%$ of its proposals yet is worth only $14$ independent
draws. Explain how both facts can be true at once.
\answerspace{4em}

\item Why is a burn-in discarded at all, and which of the three runs would suffer most from
\emph{not} discarding it?
\answerspace{4em}
\end{parts}

%%%%%%%%%%%%%%%%%%%%%%%%%%%%%%%%%%%%%%
\item \textbf{The three numbers.} Use posterior sd $= 0.586$ throughout.

\begin{parts}
\item Compute the MCSE of the posterior mean for the good run (ESS $867$) and for the bad one
(ESS $14.2$).
\answerspace{4em}

\item You want the MCSE of the good run cut in half. What ESS do you need, and roughly how much
longer must you run the chain?
\answerspace{3.5em}

\item Four chains from spread-out starting values give $\widehat{R} = 1.574$ at a step of $0.05$, and
$\widehat{R} = 1.001$ at a step of $1.20$. Say what $\widehat{R}$ compares, and what the first value
tells you that ESS alone would not.
\answerspace{4.5em}

\item Which of the three numbers answers ``have I converged?'' and which answer ``have I run long
enough?''
\answerspace{3.5em}
\end{parts}

%%%%%%%%%%%%%%%%%%%%%%%%%%%%%%%%%%%%%%
\item \textbf{Everything by counting.} The good run leaves $5000$ retained draws of $\lambda$. Of
them, $1570$ exceed $3$, the $2.5$th percentile is $1.72$, and the $97.5$th is $4.01$.

\begin{parts}
\item Give the point estimate, the $95\%$ credible interval, and $P(\lambda > 3)$, saying for each
what you did to the draws.
\answerspace{4.5em}

\item Compare all three with the exact values in the scenario. Are the approximations good enough to
report?
\answerspace{3.5em}

\item Write the one-sentence interpretation of the credible interval. Then say what makes it a
different claim from a frequentist confidence interval with the same endpoints.
\answerspace{4.5em}

\item The station manager asks: ``Can we reject the claim that we average at most $3$ sightings a
night, at the $5\%$ level?'' Answer the question a Bayesian can actually answer, and say what is
missing from hers.
\answerspace{4.5em}
\end{parts}

%%%%%%%%%%%%%%%%%%%%%%%%%%%%%%%%%%%%%%
\item \textbf{Predicting tomorrow night, and find the bugs.} Two predictions for the count on one
new night:

\begin{center}
\renewcommand{\arraystretch}{1.2}
\begin{tabular}{lccc}
\toprule
 & $P(Y \ge 5)$ & \textbf{SD} & \textbf{95\% interval} \\
\midrule
Plug-in: Poisson$(2.750)$ & 0.145 & 1.658 & $[\,0,\; 6\,]$ \\
Posterior predictive & 0.154 & 1.768 & $[\,0,\; 7\,]$ \\
\bottomrule
\end{tabular}
\end{center}

\begin{parts}
\item Explain in one sentence why the predictive column is the wider one. Name the source of
uncertainty the plug-in row leaves out.
\answerspace{3.5em}

\item Describe the recipe that produced the predictive row, in terms of the $5000$ draws.
\answerspace{3.5em}

\item The code below is meant to fit and predict. It runs, prints plausible numbers, and is wrong in
\textbf{three} separate ways.

\begin{lstlisting}[language=R]
chain  <- rwmh(step = 0.05, n = 6000)   # acceptance 0.987, looks great
lambda <- chain$x

mean(lambda)
quantile(lambda, c(0.025, 0.975))

rpois(5000, mean(lambda))
\end{lstlisting}

Bug 1 is the step size chosen, bug 2 is what never happens to \texttt{lambda}, and bug 3 is the last
line. Name all three and give the corrections.
\answerspace{5.5em}

\item The same classmate's write-up ends with the two sentences below. Mark each \textbf{defensible}
or \textbf{not}, and rewrite the ones that are not.
\begin{enumerate}[label=\textbf{S\arabic*.}, leftmargin=2.2em, itemsep=.25em, topsep=.35em]
\item ``Our sampler accepted $98.7\%$ of proposals, so it explored the posterior very well.''
\item ``The posterior probability that $\lambda > 3$ is $0.31$, so we reject $\lambda \le 3$ at the
$5\%$ level.''
\end{enumerate}
\answerspace{5em}
\end{parts}

\end{questions}

\end{document}
